\documentclass{article}
\usepackage{pathfinder-paper}

\title{An Information Limit for Pairwise Quote Judgments in a Double Auction}

\begin{document}
\maketitle

\begin{abstract}
We examine whether a pairwise judge given an acting trader's permitted information can rank two quotes by
their counterfactual final surplus in a private-value double auction. In a constructed continuation of the
auction studied by Struski and colleagues, two worlds preserve the full reservation-value schedule and the
trader's information, yet reverse the welfare ranking of a crossing bid and an improving bid. With equal
conditional probability on those worlds, any such judge has at most one-half accuracy on the realised
better-quote label. The optimal local decision still has a well-defined expected-welfare ranking; the
specified value of learning a competing buyer's valuation is \(0.25\) surplus units. This is a conditional
information boundary for an adaptation of Wang and colleagues' pairwise credit method, not an empirical
estimate of auction frequency or a guarantee about market intervention.
% ledger: 11, 12, 14, 15, 19, 20, 22, 23, 28
\end{abstract}

\section{Question and scope}

Q, \emph{Competitive Market Behavior of LLMs}, studies prompted, privately informed traders in a continuous
double auction and reports slower or absent convergence and missed gains from trade relative to a human
benchmark \cite{struski2026competitive}. Its allocative-efficiency measure is realised buyer-plus-seller
surplus divided by maximum available surplus. A profitable spread crossing can recover a missed trade, but it
can also allocate a seller's unit to a lower-value buyer before a higher-value buyer arrives.
% ledger: 1, 9, 17

P, \emph{MARS-RA}, compares \emph{agents} at a state using multimodal-model judgments, aggregates the
comparisons into Bradley--Terry scores, and uses those scores in potential-based reward shaping for
cooperative learning \cite{wang2026marsra}. Its stated statistical target is the comparator's latent
preference under its assumptions; economic welfare accuracy is not established by that result. Comparing
\emph{alternative quotes of one trader} is therefore an adaptation of P's judgment idea, not an output of P's
agent scores. The question here is whether that adapted judge can identify the quote with greater final market
surplus when its input is restricted to Q's acting-trader information.
% ledger: 1, 10, 12, 17, 19, 22, 23

Q has eleven buyers and eleven sellers, each able to trade once per round. Reservation values on each side
occupy the \(0.75\)-to-\(3.25\) grid in increments of \(0.25\). A round lasts at most 300 iterations. At each
iteration Q selects an agent uniformly from those who have not traded. A buyer's bid at or above the standing
ask executes at that ask; a bid above the standing bid and below the ask becomes the new bid. The trader sees
its own reservation value, standing quotes, public offer history, and its own private history, but not other
traders' reservation values. Q instructs traders to maximise private profit.
% ledger: 7, 9, 17, 19, 20, 22, 28

\section{A conditional information boundary}

Let \(W(a,\omega)\) denote final realised surplus after quote \(a\) under the specified continuation in world
\(\omega\). The comparison concerns a crossing bid \(a_{\mathrm{c}}=1.50\) and an improving bid
\(a_{\mathrm{i}}=1.49\), both available to the same buyer. It does not compare agents or infer a quote ranking
from P's Bradley--Terry scores.
% ledger: 19, 20, 22, 23

\begin{proposition}[Identical local information, opposite quote rankings]
Consider the following constructed states, with the acting buyer \(B_2\) selected just before one final
iteration. A seller \(S\) has cost \(0.75\) and the sole relevant standing ask \(1.50\); \(B_2\) has value
\(2.25\), and the standing bid is below \(1.49\). Buyers \(B_1\) and \(B_3\) have neither posted nor otherwise
revealed their values. In world A, their values are \(3.25\) and \(1.75\), respectively; in world B, those
values are swapped. Their identities and the common history convey no signal of the assignment to \(B_2\).
Couple the continuation so that \(B_1\) is selected on the final turn in either branch and bids \(1.50\) if
\(S\) remains; \(B_3\) does not act, and all other welfare contributions coincide. Then \(B_2\)'s permitted
information and Q's full buyer-value multiset are identical across worlds, while the sign of
\(W(a_{\mathrm{c}},\omega)-W(a_{\mathrm{i}},\omega)\) is opposite.
% ledger: 11, 19, 20, 22, 25, 28
\end{proposition}

\begin{proof}
Both bids are legal under Q's quote rule. Crossing matches \(B_2\) and \(S\) at once, yielding
\(2.25-0.75=1.50\) surplus in either world. Improving leaves \(S\) available for \(B_1\)'s final crossing,
yielding \(3.25-0.75=2.50\) in A and \(1.75-0.75=1.00\) in B. Common welfare contributions cancel from the
difference. Hence the crossing-minus-improving differences are \(-1.00\) and \(+0.50\). Swapping the two
unseen values preserves the fixed value schedule, and the stipulated histories make the two observations
equal.
% ledger: 19, 20, 22, 25
\end{proof}

Under an equal conditional prior on A and B, a judge measurable only from that common information, including
one that repeats the same pairwise comparison, has at most \(1/2\) accuracy on the \emph{realised}
better-quote label. This is an information ceiling for the stated two-world distribution, not a measured error
rate in Q's simulations. A view that includes \(B_1\)'s private value can distinguish the worlds; access to
every market valuation is sufficient but unnecessary for this example.
% ledger: 11, 12, 13, 20, 22

The ceiling does not imply an inability to make the best local decision. Improving has conditional expected
surplus \((2.50+1.00)/2=1.75\), compared with \(1.50\) from crossing. A decision maker who sees \(B_1\)'s
value improves in A and crosses in B, obtaining \((2.50+1.50)/2=2.00\) in expectation. The \(0.25\) difference
is the value of that hidden information for this prior and continuation. Repeated judgments can reduce
sampling noise in a comparator, but cannot remove this information gap.
% ledger: 12, 14, 15, 20

\section{Relation to existing methods}

The action-rollout step is already present in Chen and colleagues' cooperative-LLM credit method: they restore
decision points, vary actions, and evaluate continuations under a fixed behaviour policy \cite{chen2026exact}.
The example here supplies a different target and constraint, namely final auction surplus under private
reservation values. Jadhav and colleagues have already used an LLM critic to shape bidding in a partially
observed continuous double auction for electricity trading \cite{jadhav2025fairness}; an LLM-guided auction
intervention is therefore not a contribution of this paper. Neither of these works, as examined here,
establishes the stated fixed-schedule, two-world quote-ranking calculation.
% ledger: 7, 16, 17, 23

For a fixed matched buyer and seller with value \(v\), cost \(c\), and transaction price \(p\), joint profit satisfies
\[
(v-p)+(p-c)=v-c.
\]
Thus an allocation-fixed price change transfers rent without changing that trade's joint surplus. This
supports a control for whether a judge mistakes rent capture for welfare creation. If a changed quote changes
subsequent behaviour or the allocation, the continuation must instead determine its total welfare effect.
% ledger: 3, 23, 26

If the comparison expands beyond two quotes, P's unregularised Bradley--Terry fit has a relevant boundary:
\(K\) unanimous comparisons favouring one of two entities give a connected undirected graph but loss
\(L(d)=K\log(1+e^{-d})\), whose infimum has no finite score difference \(d\). A penalised fit or symmetric
pseudocounts would address this separation in a larger quote-set study. The present two-quote result needs no
Bradley--Terry fit.
% ledger: 4, 5, 6, 23, 26

P's potential-based shaping also has a distinct scope. With fixed initial and zero terminal potential, its
discounted shaping sum for agent \(i\) telescopes to
\[
\sum_{t=0}^{T-1}\gamma^t\bigl(\gamma\Phi_i(s_{t+1})-\Phi_i(s_t)\bigr)=-\Phi_i(s_0).
\]
Under these conditions it cannot by itself change a strategic equilibrium. Q uses prompted profit-maximising
traders rather than policies trained on P's cooperative reward. An agent's exit after trade requires potential
handling on that transition for the identity to apply; the record does not show a violation in P's
implementation.
% ledger: 7, 8, 20

\section{Interpretation}

The example separates a realised better-quote label from the best decision given the available information. In
the specified balanced prior, a locally informed decision maker can choose improving and achieve its
conditional expected-welfare optimum while disagreeing with half of the realised-world labels.
% ledger: 12, 14, 15, 20

\section{Limitations}

The proposition is an existence result conditional on the specified final selection, bidding continuation, and
absence of identity leakage. It does not establish that these states arise often, or that Q's prompted traders
would choose the stipulated bid. Because a trade changes the set of eligible agents, counterfactual replays
should couple underlying random draws and reapply Q's selection rule in each branch, rather than impose an
identical realised actor sequence. P already acknowledges that tasks depending on unobserved domain-specific
state may fall outside its scope. The calculation quantifies that boundary in this market; it does not refute
P's stated limitation.
% ledger: 19, 20, 22, 23, 25, 28

\section{Open questions}

A direct test would locate Q states where both quotes are legal, replay them with fixed continuation policies
and coupled randomness, and label the difference in final surplus. It would measure state frequency and
compare local-information pairwise judgments with direct conditional-welfare prediction and a judge given the
relevant hidden value. These are proposed measurements, not findings. The effect size in Q, any predictive
advantage of pairwise judgment, rollout feasibility, and any downstream gain in market efficiency remain open.
% ledger: 17, 19, 23, 25, 28

\bibliographystyle{plainurl}
\bibliography{references}
\end{document}
